% =============================================================================
\chapter*{Abstract}
\addcontentsline{toc}{chapter}{Abstract}

\nota{Versão equivalente ao Resumo, expandindo a partir do abstract finalizado para o Congresso de Genética 2026 (2.485 caracteres). Manter mesma estrutura seccional.}

\reaproveitar{Base do abstract congressual (adaptar tempo verbal e expandir para 250-500 palavras):}

\begin{quote}
\textit{Chemical contamination by priority pollutants remains a global concern, reflected in international regulatory frameworks. Functional, toxicological and regulatory evidence remain scattered across independent databases, hindering multi-omics analysis of these compounds. \biorempp{} (Bioremediation Potential Profile) addresses this gap by unifying regulatory prioritisation with functional annotation through a compound-centric approach, delivered as an ecosystem of four components sharing a curated evidence base consolidated in the \biorempp{} core database: the Integrated Database, the Database Explorer (DBX), the analysis-based Web Service (BWS) and a command line interface (CLI). [\ldots]}
\end{quote}

\escrever{Expandir os blocos de metodologia e resultados no abstract para tese; adicionar frase sobre validação dupla e sobre as três publicações de aplicação retrospectiva (Freitas 2023, Silva-Portela 2024, Freitas 2025).}

\vspace{1em}
\textbf{Keywords:} bioinformatics; bioremediation; functional genomics; multi-omics integration; priority compounds; tool ecosystem.

